\documentclass[12pt, a4paper]{article}
\usepackage[utf8]{inputenc}
\usepackage[italian]{babel}
\usepackage{amsmath, amssymb, amsthm}
\usepackage{graphicx}
\usepackage{geometry}
\geometry{a4paper, top=3cm, bottom=3cm, left=2.5cm, right=2.5cm}
\usepackage{fancyhdr}
\usepackage{color}
\usepackage{hyperref}

% Definizione per Teoremi e Definizioni per matchare lo stile del PDF
\newtheorem{definition}{Definition}[section]
\newtheorem{theorem}{Law}[section]

% Setup dell'intestazione (Header e Footer) identica al PDF
\pagestyle{fancy}
\fancyhf{}
\fancyhead[L]{\nouppercase{\leftmark}}
\fancyhead[R]{Bandi Tommaso Felice}
\fancyfoot[C]{\thepage}

\begin{document}

% ------ PAGINA DEL TITOLO ------
\begin{titlepage}
    \centering
    \vspace*{4cm}
    {\Huge \bfseries ELETTROTECNICA \par}
    \vspace{1.5cm}
    {\Large Bandi Tommaso Felice \par}
    \vspace{0.5cm}
    {\large 5 agosto 2023 \par}
\end{titlepage}

% ------ INDICE ------
\tableofcontents
\newpage

% ------ CAPITOLO 1 ------
\section{Introduzione al Corso ELETTROTECNICA}

\subsection{Obiettivi dell'insegnamento}
Il corso si propone di fornire agli studenti nozioni fondamentali di teoria dei circuiti e cenni di elettromagnetismo. Con l'obiettivo di far acquisire loro l'abilità di applicare le nozioni teoriche per la risoluzione dei problemi, in ambito dell'ingegneria dell'informazione e industriale, attraverso l'analisi ed i metodi risolutivi. Vengono inoltre fornite le metodologie deduttive dell'analisi matematica e fisica applicati a casi pratici dell'ingegneria.

\subsection{Risultati di apprendimento attesi}
\begin{itemize}
    \item Conoscere cos'è un circuito elettrico ed i suoi componenti fondamentali.
    \item Conoscere il concetto di modello circuitale e i suoi limiti di applicabilità.
    \item Conoscere le leggi ed i teoremi per lo studio dei circuiti fondamentali (e limiti di validità), necessari per affrontare l'analisi di un circuito.
    \item Comprendere il campo elettromagnetico nei circuiti elementari.
    \item Comprendere il legame tra le leggi dei campi elettro-magnetici ed il modello circuitale nonché le principali applicazioni dei campi stessi.
    \item Saper utilizzare le conoscenze acquisite per formulare e risolvere il problema dell'analisi di circuiti elettrici.
    \item Essere in grado di interpretare in modo critico i risultati ottenuti.
\end{itemize}

\subsection{Argomenti trattati}
\begin{enumerate}
    \item \textbf{Introduzione. 1. Circuiti elettrici: cosa modella il fenomeno fisico.} 2. Il concetto di bipolo. 3. Le grandezze elettriche: tensione, corrente e potenza. 4. Unità di misura. Voltmetro e amperometro. 5. Leggi di Kirchhoff delle correnti e delle tensioni. 6. Potenza ed energia. Teorema di Tellegen.
    
    \item \textbf{Bipoli resistivi e circuiti adinamici.} 1. Il resistore ideale. 2. Bipoli attivi ed elementi passivi, generatori ideali di tensione e di corrente, corto circuito e circuito aperto. 3. Modelli di Thevenin (serie) e di Norton (parallelo). 
    
    \item \textbf{Bipoli lineari a più porte.} 1. Matrice di resistenza $R$, di conduttanza $G$, ed ibrida $H, H'$. 2. Potenza in un n-polo lineare. 3. Trasformatori ideali.
    
    \item \textbf{Analisi dei circuiti.} 1. Topologia di una rete. 2. Analisi nodale di circuiti (metodo dei potenziali ai nodi). 3. Modelli modificati. 4. Teorema della sovrapposizione degli effetti.
\end{enumerate}

\subsection{Prerequisiti}
Matematica: Calcolo differenziale e integrale per funzioni di variabile reale. Algebra dei numeri complessi. Elementi di calcolo matriciale.\\
Fisica: Potenza, Lavoro, energia.

\newpage

% ------ CAPITOLO 2 ------
\section{Introduzione}
\subsection{Vettori}
\subsubsection{Definizione e norma euclidea}

\begin{definition}[Vettori]
Un vettore è una grandezza matematica multi-direzionale di norma, direzione e verso, può essere un vettore colonna o un vettore riga.
\end{definition}

In fisica un vettore descrive la posizione, quindi le coordinate in un punto in uno spazio n-dimensionale.
La norma (norma euclidea, norma 2) di un vettore è la lunghezza del vettore stesso:
\[
||v|| = ||v||_2 = \sqrt{\sum_{i=1}^n v_i^2}
\]

\subsubsection{Prodotto scalare-vettore}
Il prodotto tra uno scalare e un vettore dà come risultato un vettore multiplo del vettore primario:
\[
\alpha \cdot \vec{v} = \begin{bmatrix} \alpha v_1 \\ \alpha v_2 \\ \vdots \\ \alpha v_n \end{bmatrix}
\]

\subsubsection{Somma e differenza}
La somma tra due vettori corrisponde a sommare le singole coordinate dei due vettori e graficamente può essere vista traslando il secondo vettore sulla punta del primo.
La differenza tra due vettori corrisponde a sommare al primo vettore l'opposto del secondo. Graficamente si traccia partendo dalla punta del secondo vettore all'estremità del primo.

\subsubsection{Versori e normalizzazione}
\begin{definition}[Versore]
Un versore è un vettore con norma unitaria:
\[
||v|| = 1
\]
\end{definition}
Per normalizzare un vettore si divide il vettore per la sua norma.

\subsubsection{Prodotto scalare (dot product)}
Il prodotto scalare tra due vettori dà come risultato la somma del prodotto tra le singole componenti corrispondenti dei vettori:
\[
\vec{v} \cdot \vec{u} = \sum_{i=1}^n v_i u_i
\]

\subsubsection{Prodotto vettoriale}
\begin{definition}[Prodotto vettoriale]
Siano $\vec{u}, \vec{v}$ due vettori, allora il prodotto vettore $\vec{w} = \vec{u} \times \vec{v}$ ha:
\begin{itemize}
    \item Modulo: $||\vec{u} \times \vec{v}|| = ||u|| \cdot ||v|| \cdot \sin(\theta)$
    \item Direzione: perpendicolare al piano su cui giacciono i due vettori
    \item Verso: segue la regola della mano destra.
\end{itemize}
\end{definition}

\newpage

% ------ CAPITOLO 3 ------
\section{Elettrostatica}
\subsection{Carica e campo}
\subsubsection{Definizione di elettromagnetismo}

\begin{definition}[Elettromagnetismo]
La branca della fisica classica che studia le cariche elettriche e i campi elettrici da esse generati.
\end{definition}

La carica è una caratteristica fondamentale della materia (ad esempio come la massa) ed è quantizzata, ciò significa che si trova in multipli interi della carica fondamentale dell'elettrone:
\[
e = -1.6 \cdot 10^{-19} \, \text{C}
\]
Per le dimensioni dei problemi è considerata buona come ipotesi che stiamo operando nel vuoto e le cariche sono considerabili puntiformi.

\subsubsection{Legge di Coulomb}
È fondamentale per studiare l'interazione tra le cariche puntiformi:
\[
\vec{F} = \frac{1}{4\pi\varepsilon_0} \frac{q_1 q_2}{r^2} \hat{u}_r
\]
In questa formula compare la costante dielettrica del vuoto che ha valore: $\varepsilon_0 = 8.854 \cdot 10^{-12} \, \text{F/m}$.
Inoltre si dimostra come due cariche dello stesso segno si respingono e due cariche di segno opposto si attraggono.

% NOTA PER L'UTENTE: Le immagini disegnate a mano non possono essere generate nativamente dal codice di testo LaTeX. 
% Devi fare uno screenshot del disegno e salvarlo (es. come "coulomb.png"), poi de-commentare le righe qui sotto per inserirlo.
\begin{figure}[h]
    \centering
    % \includegraphics[width=0.6\textwidth]{coulomb.png}
    \caption{Legge di Coulomb}
\end{figure}

\subsubsection{Campo elettrico}
\begin{definition}[Campo]
Definisce l'influenza esercitata da una sorgente su uno spazio. Ci possono essere due tipologie di campo:
\begin{itemize}
    \item Campo scalare (es. Campo di temperature)
    \item Campo vettoriale (es. Campo gravitazionale, elettrico)
\end{itemize}
\end{definition}

\subsubsection{Linee di campo}
\begin{definition}[Linee di campo]
Le linee di campo sono l'insieme dei punti in cui un dato campo vettoriale è tangente.
\end{definition}

\begin{definition}[Campo Elettrico]
Supponiamo di avere una carica $q$ che genera un campo elettrico $\vec{E}$.
\[
\vec{E} = \frac{\vec{F}}{q_0}
\]
\end{definition}

\begin{definition}[Carica di prova]
Un corpo puntiforme il cui campo può essere considerato esente di carica di prova se:
\begin{itemize}
    \item La carica è soggetta alla forza generata dalle altre cariche;
    \item La carica non altera il campo elettrico totale.
\end{itemize}
\end{definition}

\begin{definition}[Campo Elettrico]
Supponiamo di avere una carica $Q$ che genera un campo elettrico $\vec{E}$.
\[
\vec{E} = \frac{\vec{F}}{q_0}
\]
Indica il campo elettrico e dice se:
\begin{itemize}
    \item Stazionario: il campo non varia nel tempo, quindi la derivata parziale rispetto al tempo è nulla:
    \[
    \frac{\partial \vec{E}}{\partial t} = 0
    \]
    \item Quasi stazionario: il campo varia nel tempo molto lentamente rispetto alle variazioni interne al circuito, quindi la derivata parziale rispetto al tempo è molto prossima allo zero:
    \[
    \frac{\partial \vec{E}}{\partial t} \simeq 0
    \]
\end{itemize}
\end{definition}

Cosa genera un campo elettrico??
Siccome vale la Legge di Coulomb:
\[
\vec{E} = \frac{\vec{F}}{q_0} = \frac{1}{4\pi\varepsilon_0} \frac{Q}{r^2} \hat{u}_r
\]
Quindi il campo elettrico generato da una carica puntiforme è diretto radialmente rispetto alla carica stessa e il verso dipende dal segno della carica:
\begin{itemize}
    \item Carica positiva $\implies$ radiale e uscente
    \item Carica negativa $\implies$ radiale ed entrante
\end{itemize}

\begin{figure}[h]
    \centering
    % \includegraphics[width=0.8\textwidth]{linee_di_campo.png}
    \caption{Linee di campo, lineari di flusso}
\end{figure}

\newpage

% ------ SEZIONE 3.2 ------
\subsection{Tensione, flusso e corrente}
\subsubsection{Tensione elettrica}
\begin{itemize}
    \item Poniamo la carica di prova in una regione di spazio in cui sia presente un campo elettrico $\vec{E}$
    \item La carica di prova andrà a spostarsi lungo le linee di campo del campo elettrico $\vec{E}$
    \item Il campo elettrico compie un lavoro sulla carica di prova $q$
    \item Vogliamo calcolare il lavoro necessario per spostare la carica di prova $q$ in presenza del campo elettrico $\vec{E}$ lungo un generico percorso $\gamma$ (da $A$ a $B$).
    \item Vogliamo spostarci da $A$ a $B$ lungo $\gamma$ e consideriamo varie spezzate abbastanza piccole da approssimare con una buona precisione la curva $\gamma$.
\end{itemize}

\begin{figure}[h]
    \centering
    % \includegraphics[width=0.6\textwidth]{spostamento_carica.png}
    \caption{Spostamento carica su linea curva}
\end{figure}

Il segno meno è da interpretare come lavoro in opposizione al campo elettrico, infatti questo è il lavoro che compiamo noi, non quello relativo allo spostamento della carica lungo la curva.
\[
L_{AB} = \int_A^B \vec{F} \cdot d\vec{s} = \int_A^B q \vec{E} \cdot d\vec{s} = q \int_A^B \vec{E} \cdot d\vec{s}
\]

Infine pensando al limite di lunghezze di spostamento molto piccole $\Delta s \to 0$ (limite che vedremo in Analisi 2)

\[
\Delta L = - \vec{F} \cdot d\vec{s} = - q \vec{E} \cdot d\vec{s}
\]

\begin{definition}[Tensione elettrica]
Lungo una qualunque linea chiusa, la somma algebrica delle tensioni, prese con il segno opportuno in base al verso di percorrenza della maglia, è nulla.
Valgono le relazioni:
\[
\oint_\gamma \vec{E} \cdot d\vec{s} = 0
\]
\end{definition}

Quindi segue che:
\[
V_{AB} + V_{BA} = 0 \implies V_{AB} = -V_{BA}
\]

Dimostrazione della legge di Kirchhoff per le tensioni: è una generalizzazione del calcolo della tensione elettrica lungo un percorso da A a B. Il lavoro svolto per spostare la carica da A a B è:
\[
V_{AB} = \frac{L_{AB}}{q} = \int_A^B \vec{E} \cdot d\vec{s}
\]

\newpage

\begin{figure}[h]
    \centering
    % \includegraphics[width=0.6\textwidth]{kirchhoff_tensioni.png}
    \caption{Legge di Kirchhoff per le tensioni}
\end{figure}

Siccome il prodotto scalare:
\[
\vec{E} \cdot d\vec{s} = |\vec{E}| \cdot |d\vec{s}| \cos \theta
\]
Ora calcoliamo il lavoro sfruttando il fatto che possiamo portar fuori dall'integrale le costanti e che possiamo parametrizzare la curva in funzione di $t$:
\[
L_{AB} = \int_A^B -q \vec{E} \cdot d\vec{s} = -q \int_A^B \vec{E} \cdot d\vec{s} = -q \int_{t_A}^{t_B} \vec{E} \cdot \frac{d\vec{s}}{dt} dt = -q \int_{t_A}^{t_B} \vec{E} \cdot \vec{v} dt = W_B - W_A = -q(V_B - V_A)
\]
Dove abbiamo che:
\begin{itemize}
    \item $W$ è l'energia potenziale
    \item $V_{A}$ e $V_{B}$ sono i \textit{potenziali elettrici} nei punti A e B
    \item $V_{AB}$ è la \textit{differenza di potenziale}
\end{itemize}

Inoltre da questo si ricava che:
\[
V_{AB} = \int_A^B \vec{E} \cdot d\vec{s}
\]

Allora il lavoro non dipende dal percorso, ma solo dai punti A e B, quindi possiamo dire che il campo elettrico è un campo conservativo.
Possiamo estendere questa dimostrazione anche al caso dei campi elettrici quasi stazionari. Consideriamo i due percorsi $\gamma_1$ e $\gamma_2$:
Allora il lavoro rispettivo sarà:
\[
L_{\gamma_1} = -q (V_B - V_A), \quad L_{\gamma_2} = -q (V_A - V_B)
\]
Ora considero il percorso chiuso $\gamma = \gamma_1 + \gamma_2$ allora il lavoro totale su $\gamma$ è:
\[
L_\gamma = L_{\gamma_1} + L_{\gamma_2} = 0
\]

\newpage

Quindi l'integrale di linea, su una linea chiusa (detto circuitazione) è:
\[
\oint_\gamma \vec{E} \cdot d\vec{s} = 0
\]
Di conseguenza la circuitazione e la tensione complessiva del campo elettrico lungo una linea chiusa sono nulle.

\subsubsection{Flusso di un campo vettoriale}
\begin{definition}[Flusso]
Definiamo flusso di un campo $\vec{A}$ attraverso la superficie $S$ come: $A \cdot \hat{u}_n \, dS$
\end{definition}

\begin{figure}[h]
    \centering
    % \includegraphics[width=0.6\textwidth]{flusso.png}
\end{figure}

Approssimiamo il flusso $\Phi$ attraverso una superficie estesa $S$ come:
\[
\Phi(\vec{A}) = \int_S \vec{A} \cdot \hat{u}_n \, dS
\]
Dove $d\vec{S}$ sta a significare che in ogni piccola quadratino il campo elettrico può essere assunto come costante.
Prendo il limite per $dS \to 0$ allora:
\[
\Phi(\vec{A}) = \int_S \vec{A} \cdot d\vec{S}
\]

Calcoliamo il flusso di $\vec{E}$ generato da una carica puntiforme attraverso una superficie sferica. Abbiamo che il raggio della sfera è $r$, e troviamo che il campo è:
\[
\vec{E} \cdot \hat{u}_n = |\vec{E}| = \frac{1}{4\pi\varepsilon_0} \frac{q}{r^2}
\]

Facendo i conti:
\[
\Phi(\vec{E}) = \oint_S \vec{E} \cdot d\vec{S} = \oint_S |\vec{E}| \, dS = \frac{1}{4\pi\varepsilon_0} \frac{q}{r^2} \oint_S dS = \frac{1}{4\pi\varepsilon_0} \frac{q}{r^2} 4\pi r^2 = \frac{q}{\varepsilon_0}
\]

\newpage

\begin{theorem}[Legge di Gauss per il flusso del campo elettrico]
\[
\Phi(\vec{E}) = \frac{Q_{\text{int}}}{\varepsilon_0}
\]
\end{theorem}
Possiamo generalizzare il risultato appena ottenuto a qualsiasi superficie in cui sono racchiuse le cariche.
Dimostrazione aggiuntiva:

\begin{figure}[h]
    \centering
    % \includegraphics[width=0.8\textwidth]{gauss_proof.png}
\end{figure}

\[
\Phi(\vec{E}) = \oint_S \vec{E} \cdot \hat{u}_n \, dS = \oint_S \vec{E} \cdot d\vec{S} = \frac{q}{\varepsilon_0}
\]

\[
d\vec{S} = \hat{u}_n \, dS, \quad \vec{E} = \frac{q}{4\pi\varepsilon_0 r^2} \hat{u}_r
\]

\[
\vec{E} \cdot d\vec{S} = \frac{q}{4\pi\varepsilon_0 r^2} \hat{u}_r \cdot \hat{u}_n \, dS = \frac{q}{4\pi\varepsilon_0 r^2} \cos\theta \, dS
\]

Il flusso totale attraverso la sfera coincidente con la superficie chiusa:
\[
\oint_S \vec{E} \cdot d\vec{S} = \oint_S \frac{q}{4\pi\varepsilon_0 r^2} \cos\theta \, dS = \frac{q}{\varepsilon_0}
\]

\newpage

\subsubsection{Corrente elettrica}
\begin{definition}[Corrente elettrica]
Definiamo corrente elettrica la quantità di carica che attraversa la superficie $S$ in un intervallo di tempo $\Delta t$:
\[
I = \frac{\Delta q}{\Delta t}
\]
\end{definition}

Passando al limite per istanti infinitesimali:
\[
i(t) = \lim_{\Delta t \to 0} \frac{\Delta q}{\Delta t} = \frac{dq}{dt}
\]

La corrente quindi è una quantità scalare e ha unità di misura $[A] = [\text{C/s}]$.
Ora diamo una definizione più rigorosa di corrente, consideriamo la seguente figura.

\begin{figure}[h]
    \centering
    % \includegraphics[width=0.5\textwidth]{corrente_superficie.png}
    \caption{Carica che attraversa una superficie S}
\end{figure}

La Quantità di carica che attraversa la superficie $S$ in un intervallo di tempo $dt$, dove $\hat{u}_n$ è la normale alla superficie, vale:
\[
dq = \rho \, dV = \rho (\vec{v} \cdot \hat{u}_n \, dt) \, dS
\]

Abbiamo che il vettore $\vec{J} = \rho \vec{v}$ è detto \textit{densità di corrente}.

\begin{definition}[Densità di corrente]
La corrente $I$ è il flusso del vettore densità di corrente $\vec{J}$ attraverso la superficie $S$:
\[
I = \int_S \vec{J} \cdot \hat{u}_n \, dS
\]
L'unità di misura di $J$ è $[\text{A/m}^2]$.
\end{definition}

\end{document}
